% called by main.tex
%
\chapter{Despliegue con Docker desde GitHub}
\label{anexo::docker-ghcr}

% [PENDIENTE: verificar el nombre y la etiqueta exactos de la imagen publicada y el
%  contenido de docker-compose.ghcr.yml antes de la entrega.]

A diferencia del Anexo~\ref{anexo::docker-local}, que construye la imagen a partir del
código fuente, este anexo describe cómo ejecutar el sistema descargando directamente la
\textbf{imagen ya publicada} en GitHub Container Registry (GHCR), sin necesidad de
compilarla.

\section{Descarga de la imagen}
\label{anexo::ghcr_pull}

La imagen del proyecto está publicada en GHCR y puede descargarse con:

\begin{verbatim}
docker pull ghcr.io/skullsupernova/mdl_tfm:latest
\end{verbatim}

La publicación se realiza de forma automática mediante un flujo de trabajo de GitHub
Actions (\texttt{.github/workflows/docker-publish.yml}) en cada actualización del
repositorio.

\section{Ejecución}
\label{anexo::ghcr_run}

El repositorio incluye un fichero de orquestación específico
(\texttt{docker-compose.ghcr.yml}) que utiliza la imagen prepublicada en lugar de
construirla localmente. Desde la raíz del repositorio:

\begin{verbatim}
docker compose -f docker-compose.ghcr.yml up
\end{verbatim}

Los servicios quedan accesibles en las mismas direcciones que en el despliegue local
(API en \texttt{:8000}, interfaz web en \texttt{:8501}), con resultados idénticos. Esta
vía es la más rápida para evaluar el sistema, al evitar la construcción de la imagen.
